\documentclass[11pt]{article}
\usepackage[a4paper,margin=30mm]{geometry}
\usepackage{amsmath,amssymb,amsthm,booktabs,microtype}
\usepackage[hidelinks]{hyperref}

\newtheorem{theorem}{Theorem}
\newtheorem{lemma}[theorem]{Lemma}
\newtheorem{corollary}[theorem]{Corollary}
\theoremstyle{remark}
\newtheorem{remark}[theorem]{Remark}

\newcommand{\Z}{\mathbb Z}
\newcommand{\Q}{\mathbb Q}
\newcommand{\vp}{v_p}

\title{A Sharp Endpoint-Denominator Theorem\\
for Residue-Class Recurrences}
\author{}
\date{August 2026}

\begin{document}
\maketitle

\begin{abstract}
We determine the exact uniform denominator of a family of inhomogeneous
recurrences which decouple into residue classes modulo an arbitrary step
$q$.  For endpoint weight $w$, the residual modulus is
\[
 M(q,w)=\prod_{p\mid q}p^{\left\lceil
  \frac{w}{q}\left(v_p(q)+\frac1{p-1}\right)\right\rceil}.
\]
The proof combines an elementary arithmetic-progression product bound with a
unique-lowest-valuation endpoint argument.  Binomial inversion transports the
classification to the original companion basis.  The parity-square theorem,
whose sharp modulus is $4$, is the case $(q,w)=(2,2)$.  A mixed step-three
specialization proves $L_n^2$-integrality, with optimal exponent, for the
second solution of Zagier's sporadic sequence $\mathbf B$; more generally its
parameter family is integral in this sense exactly when $3$ divides the
parameter.
\end{abstract}

\section{The endpoint sequence}

Write
\[
 L_n=\operatorname{lcm}(1,2,\ldots,n),\qquad L_0=1.
\]
Fix integers $q\ge2$, $w\ge1$, and $b\in\Z$.  For every prime $p\mid q$ put
\begin{equation}\label{eq:mu}
 \mu_p(q,w)=
 \left\lceil\frac{w}{q}
 \left(v_p(q)+\frac1{p-1}\right)\right\rceil,
 \qquad
 M(q,w)=\prod_{p\mid q}p^{\mu_p(q,w)}.
\end{equation}
Let $T_m=0$ for $m\le0$, and define $T_n\in\Q$ for $n\ge1$ by
\begin{equation}\label{eq:rec}
 n^wT_n=b^q(n-q+1)^wT_{n-q}+(-b)^{n-1}.
\end{equation}
Here $(-b)^0=1$, including when $b=0$.
Equivalently, for $T(z)=\sum_{n\ge0}T_nz^n$ and
$\theta=z\,d/dz$, recurrence \eqref{eq:rec} is the endpoint operator equation
\begin{equation}\label{eq:operator}
 \{\theta^w-b^qz^q(\theta+1)^w\}T(z)=\frac{z}{1+bz}.
\end{equation}

For $0\le j<n$ such that $n-j\equiv1\pmod q$, set
\begin{equation}\label{eq:root}
 Q_q(n,j)=
 \frac{\prod_{t=0}^{r-1}(j+2+tq)}
      {\prod_{t=0}^{r}(j+1+tq)},
 \qquad r=\frac{n-j-1}{q}.
\end{equation}
An empty numerator product is $1$.

\begin{theorem}[sharp step-$q$ endpoint denominator]\label{thm:main}
For every $n\ge1$,
\begin{equation}\label{eq:finite}
 T_n=b^{n-1}
 \sum_{\substack{0\le j<n\\n-j\equiv1\pmod q}}
 (-1)^jQ_q(n,j)^w.
\end{equation}
Moreover, the following are equivalent:
\begin{enumerate}
 \item $M(q,w)\mid b$;
 \item $L_n^wb^{n-1}Q_q(n,j)^w\in\Z$ for every admissible $(n,j)$;
 \item $L_n^wT_n\in\Z$ for every $n\ge0$.
\end{enumerate}
\end{theorem}

The modulus in \eqref{eq:mu} is therefore both termwise sufficient and
globally necessary.  In particular, summation cannot lower it.

\section{The finite endpoint expansion}

\begin{proof}[Proof of \eqref{eq:finite}]
The forcing at index $j+1$ contributes $(-b)^j/(j+1)^w$.  If
$n=j+1+rq$, its propagation to $n$ multiplies this by
\[
 \prod_{t=0}^{r-1}
 \frac{b^q(j+2+tq)^w}{(j+1+(t+1)q)^w}.
\]
The total power of $b$ is $j+rq=n-1$, and the remaining quotient is
$Q_q(n,j)^w$.  Summing over the admissible source indices proves the formula.
\end{proof}

\section{The arithmetic-progression bound}

\begin{lemma}\label{lem:valuation}
Let $(n,j)$ be admissible and put $r=(n-j-1)/q$.
\begin{enumerate}
 \item If $p\nmid q$, then
 \begin{equation}\label{eq:outside}
  v_p(Q_q(n,j))\ge-v_p(L_n).
 \end{equation}
 \item If $p\mid q$ and $s=v_p(q)$, then
 \begin{equation}\label{eq:inside}
  v_p(Q_q(n,j))\ge-v_p(L_n)-rs-v_p(r!).
 \end{equation}
\end{enumerate}
\end{lemma}

\begin{proof}
Suppose first that $p\nmid q$.  For any fixed $p^a$, the denominator of
$Q_q(n,j)$ contains $r+1$ terms in one invertible residue progression, while
the numerator contains $r$ terms in another.  The former residue occurs at
most $\lceil(r+1)/p^a\rceil$ times and the latter at least
$\lfloor r/p^a\rfloor$ times.  Their difference is at most one.  Only powers
$p^a\le n$ occur, so summing over $a$ proves \eqref{eq:outside}.

Now suppose $p\mid q$.  Write the denominator factors as
$d_t=j+1+tq$, $0\le t\le r$, and choose $t_*$ so that $v_p(d_{t_*})$
is maximal.  Since $d_{t_*}\le n$,
$v_p(d_{t_*})\le v_p(L_n)$.  For $t\ne t_*$, maximality and
$d_t-d_{t_*}=q(t-t_*)$ imply
\[
 v_p(d_t)\le s+v_p(|t-t_*|).
\]
Consequently
\begin{align*}
 v_p\left(\prod_{t=0}^r d_t\right)
 &\le v_p(L_n)+rs+v_p\bigl(t_*!\,(r-t_*)!\bigr)\\
 &\le v_p(L_n)+rs+v_p(r!),
\end{align*}
because $r!/[t_*!\,(r-t_*)!]$ is an integer.  Discarding the numerator only
weakens the estimate and proves \eqref{eq:inside}.
\end{proof}

\section{Sufficiency and sharpness}

\begin{proof}[Proof of Theorem~\ref{thm:main}]
We have proved the finite formula.  Assume first that $M(q,w)\mid b$.
For $p\nmid q$, \eqref{eq:outside} makes
$L_n^wQ_q(n,j)^w$ $p$-integral.  For $p\mid q$, write
$s=v_p(q)$ and $e=v_p(b)$.  Definition \eqref{eq:mu} gives
\begin{equation}\label{eq:density}
 eq\ge w\left(s+\frac1{p-1}\right).
\end{equation}
Because $n-1=j+rq\ge rq$ and
$v_p(r!)\le r/(p-1)$,
\[
 e(n-1)\ge erq
 \ge wr\left(s+\frac1{p-1}\right)
 \ge w\bigl(rs+v_p(r!)\bigr).
\]
Together with \eqref{eq:inside}, this proves the termwise assertion, and
\eqref{eq:finite} proves $L_n^wT_n\in\Z$.

It remains to prove necessity.  The case $b=0$ is immediate, so assume
$b\ne0$.  Suppose that $p\mid q$ and
$e=v_p(b)<\mu_p(q,w)$; put $s=v_p(q)$.  At $n=qR$, the admissible source
indices are $j=q\ell-1$, $1\le\ell\le R$.  Every numerator factor of
$Q_q(qR,q\ell-1)$ is $1$ modulo $p$, while its denominator is
\[
 q^{R-\ell+1}\frac{R!}{(\ell-1)!}.
\]
The valuation of this contribution to \eqref{eq:finite} is therefore
\begin{equation}\label{eq:ellval}
 e(qR-1)-w\left((R-\ell+1)s
 +v_p\left(\frac{R!}{(\ell-1)!}\right)\right).
\end{equation}
The $\ell=1$ term is uniquely of least valuation: the difference between
\eqref{eq:ellval} and its value at $\ell=1$ is
\[
 w\bigl((\ell-1)s+v_p((\ell-1)!)\bigr)>0
 \qquad(\ell>1).
\]
Hence no cancellation is possible and
\begin{equation}\label{eq:exactval}
 v_p(T_{qR})=e(qR-1)-w\bigl(Rs+v_p(R!)\bigr).
\end{equation}
Adding $wv_p(L_{qR})$, dividing by $R$, and letting $R\to\infty$ gives
the limit
\[
 eq-w\left(s+\frac1{p-1}\right)<0.
\]
Thus $L_{qR}^wT_{qR}$ is nonintegral for all sufficiently large $R$.
Condition (3) in the theorem forces $v_p(b)\ge\mu_p(q,w)$ for every
$p\mid q$, which is exactly $M(q,w)\mid b$.
\end{proof}

\section{Binomial transport}

\begin{corollary}\label{cor:transport}
Define
\begin{equation}\label{eq:B}
 B_n=\sum_{k=0}^n\binom nk b^{n-k}T_k.
\end{equation}
Then
\[
 \boxed{\quad
 L_n^wB_n\in\Z\text{ for every }n
 \quad\Longleftrightarrow\quad M(q,w)\mid b.
 \quad}
\]
\end{corollary}

\begin{proof}
Binomial inversion gives
\[
 T_n=\sum_{k=0}^n\binom nk(-b)^{n-k}B_k.
\]
Since $L_k\mid L_n$ for $k\le n$, either transform preserves uniform
$L_n^w$-integrality.  Apply Theorem~\ref{thm:main}.
\end{proof}

At the generating-function level, \eqref{eq:B} and its inverse are exactly
\[
 B(z)=\frac1{1-bz}T\!\left(\frac{z}{1-bz}\right),
 \qquad
 T(z)=\frac1{1+bz}B\!\left(\frac{z}{1+bz}\right).
\]
Thus Corollary~\ref{cor:transport} transports the sharp arithmetic result from
the decoupled operator \eqref{eq:operator} back to its binomially conjugate
companion.

For $(q,w)=(2,2)$, one obtains $M(2,2)=4$, recovering the parity-square
endpoint theorem.  If $q=\ell$ is prime, the answer simplifies markedly:
\begin{equation}\label{eq:prime}
 M(\ell,w)=\ell^{\lceil w/(\ell-1)\rceil}.
\end{equation}
Thus $M(3,2)=3$ and $M(5,4)=5$, rather than the naive sufficient moduli
$3^2$ and $5^4$.  The parity step is exceptional in that
$M(2,w)=2^w$.

\begin{center}
\begin{tabular}{ccl}
\toprule
$q$ & $w$ & $M(q,w)$\\
\midrule
$2$ & $1,\ldots,8$ & $2,4,8,16,32,64,128,256$\\
$3$ & $1,\ldots,8$ & $3,3,9,9,27,27,81,81$\\
$4$ & $1,\ldots,8$ & $2,4,8,8,16,32,64,64$\\
$5$ & $1,\ldots,8$ & $5,5,5,5,25,25,25,25$\\
$6$ & $1,\ldots,8$ & $6,6,6,12,36,36,72,72$\\
\bottomrule
\end{tabular}
\end{center}

\section{A natural realization: Zagier's sporadic sequence
\texorpdfstring{$\mathbf B$}{B}}

The pure-power endpoint recurrence of Theorem~\ref{thm:main} is not the only
normal form with the same valuation density.  A binomial conjugation of
Zagier's quadratic recurrence~\cite{Zagier} produces a mixed step-three endpoint in which
the two missing residue classes occur once and the denominator class occurs
twice.

For $h\in\Z$, let $C^{(h)}_0=0$, $C^{(h)}_1=1$, and
\begin{equation}\label{eq:zagrec}
 (n+1)^2C^{(h)}_{n+1}
 =h(3n^2+3n+1)C^{(h)}_n-3h^2n^2C^{(h)}_{n-1}.
\end{equation}
Define
\begin{equation}\label{eq:zagtransform}
 S^{(h)}_n=\sum_{k=0}^n\binom nk(-h)^{n-k}C^{(h)}_k.
\end{equation}

\begin{theorem}[sharp step-three companion theorem]\label{thm:zagB}
For every integer $h$,
\begin{equation}\label{eq:zagclass}
 \left(\forall n\ge0,\ L_n^2C^{(h)}_n\in\Z\right)
 \quad\Longleftrightarrow\quad 3\mid h.
\end{equation}
If $h=3$, then \eqref{eq:zagrec} is the recurrence for the second solution of
Zagier's sporadic sequence $\mathbf B$.  Consequently its companion satisfies
\[
 L_n^2C^{(3)}_n\in\Z\qquad(n\ge0),
\]
and the exponent $2$ is optimal.
\end{theorem}

\begin{proof}
Let $C(x)=\sum_{n\ge0}C^{(h)}_nx^n$.  Including the initial residual,
\eqref{eq:zagrec} is equivalent to
\[
 \{\theta^2-hx(3\theta^2+3\theta+1)
        +3h^2x^2(\theta+1)^2\}C(x)=x.
\]
The generating function of \eqref{eq:zagtransform} is
\[
 S(z)=\frac1{1+hz}C\!\left(\frac{z}{1+hz}\right).
\]
Direct conjugation cancels the terms of degrees one and two and gives
\begin{equation}\label{eq:zagoperator}
 \{\theta^2+h^3z^3(\theta+1)(\theta+2)\}S(z)
 =\frac{z}{1+hz}.
\end{equation}
Thus
\begin{equation}\label{eq:zagTrec}
 n^2S^{(h)}_n+h^3(n-2)(n-1)S^{(h)}_{n-3}=(-h)^{n-1}.
\end{equation}

For $0\le j<n$ with $n-j\equiv1\pmod3$, put
\[
 r=\frac{n-j-1}{3},\qquad
 R(n,j)=
 \frac{\prod_{t=0}^{r-1}(j+2+3t)(j+3+3t)}
      {\prod_{t=0}^{r}(j+1+3t)^2}.
\]
Iteration of \eqref{eq:zagTrec} yields
\begin{equation}\label{eq:zagfinite}
 (-1)^{n-1}S^{(h)}_n
 =h^{n-1}\!
 \sum_{\substack{0\le j<n\\n-j\equiv1\pmod3}}R(n,j).
\end{equation}
Indeed, a source at $j+1$ acquires $r$ propagation signs, and
$j+r\equiv n-1\pmod2$.  In particular, for $h>0$ the right-hand side of
\eqref{eq:zagfinite} is a sum of positive rational numbers.

We prove termwise integrality when $3\mid h$.  If $p\ne3$, then for every
power $p^a$ the doubled denominator progression contains at most
$2\lceil(r+1)/p^a\rceil$ multiples of $p^a$, while the two numerator
progressions contain at least $2\lfloor r/p^a\rfloor$ in total.  The deficit
is at most two.  Hence
\begin{equation}\label{eq:zagoutside}
 v_p(R(n,j))\ge-2v_p(L_n).
\end{equation}

For $p=3$, apply the maximal-factor argument of
Lemma~\ref{lem:valuation} to
$D=\prod_{t=0}^r(j+1+3t)$.  It gives
\[
 v_3(D)\le v_3(L_n)+r+v_3(r!),
\]
and therefore
\begin{equation}\label{eq:zaginside}
 v_3(R(n,j))\ge-2v_3(L_n)-2r-2v_3(r!).
\end{equation}
If $3\mid h$, then $n-1=j+3r\ge3r$ and
$v_3(r!)\le r/2$, so
\[
 v_3\bigl(L_n^2h^{n-1}R(n,j)\bigr)
 \ge3r-2r-2v_3(r!)\ge0.
\]
Together with \eqref{eq:zagoutside}, this proves every summand of
\eqref{eq:zagfinite} integral after multiplication by $L_n^2$.  Binomial
inversion and $L_k\mid L_n$ transfer the result to $C^{(h)}_n$.

For necessity, exact propagation of \eqref{eq:zagrec} gives
\begin{equation}\label{eq:zagwitness}
 C^{(h)}_6=-\frac{39521h^5}{32400},
 \qquad
 L_6^2C^{(h)}_6=-\frac{39521h^5}{9}.
\end{equation}
Since $3\nmid39521$, integrality at $n=6$ forces $3\mid h$.  This proves
\eqref{eq:zagclass}.

For $h=3$, the first solution of \eqref{eq:zagrec} is the named sporadic row
\cite{Zagier,Gorodetsky}
\[
 A_{\mathbf B}(n)=
 \sum_{k=0}^{\lfloor n/3\rfloor}(-1)^k3^{n-3k}
 \binom n{3k}\frac{(3k)!}{k!^3},
\]
with $A_{\mathbf B}(0)=1$, $A_{\mathbf B}(1)=3$.  Hence $C^{(3)}$ is its
normalized companion.  Finally,
$C^{(3)}_2=21/4$, so $L_2C^{(3)}_2=21/2\notin\Z$; the exponent two cannot be
lowered.
\end{proof}

\section{Interpretation and reproducibility}

The exponent in \eqref{eq:mu} is a density balance.  The contribution
$v_p(q)$ is built into every step of the denominator progression, while
$1/(p-1)$ is the asymptotic valuation density of a factorial.  The factor
$1/q$ records that only one in every $q$ indices lies in a given endpoint
chain.  The unique minimum in \eqref{eq:exactval} proves that this density
cannot be improved by cancellations between endpoints.

The closest classical arithmetic-progression literature studies least common
multiples or products of consecutive progression terms; see, for example,
Hong and Qian~\cite{HongQian}.  The present statement instead determines the
sharp parameter modulus for a ratio of interlaced progressions arising from
an inhomogeneous recurrence.

The exact audit
\begin{center}
\nolinkurl{work/harmonic_jets/verify_step_q_endpoint.py}
\end{center}
checks the recurrence against \eqref{eq:finite}, the termwise bound over a
grid of $(q,w,n)$, and deficient-prime witnesses for the necessity argument.
The independent application audit
\begin{center}
\nolinkurl{work/harmonic_jets/verify_zagier_B_endpoint.py}
\end{center}
checks the operator conjugation for arbitrary $h$, the positive endpoint
formula, termwise integrality, the named binomial first row, and the two exact
sharpness witnesses in Theorem~\ref{thm:zagB}.

\begin{thebibliography}{9}
\bibitem{Gorodetsky}
O.~Gorodetsky,
\emph{New representations for all sporadic Ap\'ery-like sequences, with applications to
congruences},
Experimental Mathematics \textbf{32} (2023), no.~4, 641--656;
arXiv:2102.11839.

\bibitem{HongQian}
S.~Hong and G.~Qian,
\emph{The least common multiple of consecutive arithmetic progression terms},
arXiv:0903.0530.

\bibitem{Zagier}
D.~Zagier,
\emph{Integral solutions of Ap\'ery-like recurrence equations},
in J.~Harnad and P.~Winternitz (eds.), \emph{Groups and Symmetries: From Neolithic Scots
to John McKay}, CRM Proceedings and Lecture Notes \textbf{47}, American Mathematical
Society, 2009, 349--366.
\end{thebibliography}

\end{document}
